\scanpage{32}
\mspage{30}
infinitésimales de $\phi$, quand on se fixe celles de $G(\phi)$ à l'aide d'un cristal en groupes $\underline{G}$.

3.2. En tous cas, l'extension $(*)$ semble un invariant intéressant du groupe p-divisible envisagé. Il subsiste, par passage à la limite, en passant \deleted{\unsure{du|à}} au cas où $S$ est \deleted{un schéma formel noethérien} par exemple le spectre d'un anneau $A$ noethérien j-adique séparé et complet, où $J$ est un idéal tel que $A/J$ soit à caractéristiques résiduelles égales à $p$. On trouve par exemple un bon invariant quand $A$ est un anneau de valuation discrète complet, éventuellement d'inégales caractéristiques, à caractéristique résiduelle $p$. Dans le cas où le groupe formel\deleted{,} réduit a la structure triviale \deleted{en} mentionnée plus haut dans 3.1., et que le corps résiduel $k$ est parfait, d'où un module de Dieudonné $M$ sur $W$ (de rang $g+g'$), il semble donc, d'après la remarque faite dans 3.1., que la donnée d'un relèvement du groupe p-divisible qu'on a sur $k$ en un groupe p-divisible sur $A$, soit entièrement équivalente à la donnée d'un relèvement de la filtration de $M\otimes_Wk$ donnée par $(*)$, en une filtration de type $(g,g^*)$ de $M\otimes_WA$, \deleted{Comme} Il serait bien intéressant d'essayer de déterminer, en termes d'une telle donnée, quel est le module de Tate correspondant à la fibre générique de $\phi$ ! On aimerait préciser également, pour des groupes p-divisibles plus généraux, quelle quantité d'information est liée à l'extension $(*)$\inserted{,} qui remplace ici le $H^1$ de De Rham.
\margin{31}
